\documentclass[10pt,a4paper]{amsart}
\usepackage{amsmath,amssymb,amsthm,graphicx}
\usepackage{verbatim}
\newenvironment{reference}{\comment}{\endcomment}
%\newenvironment{reference}{}{}
\newenvironment{slogan}{\comment}{\endcomment}
\newenvironment{history}{\comment}{\endcomment}

% For commutative diagrams we use Xy-pic
\usepackage[all]{xy}

% We use 2cell for 2-commutative diagrams.
\xyoption{2cell}
% Theorem environments.
%
\theoremstyle{plain}
\newtheorem{theorem}[subsection]{Theorem}
\newtheorem{proposition}[subsection]{Proposition}
\newtheorem{lemma}[subsection]{Lemma}

\theoremstyle{definition}
\newtheorem{definition}[subsection]{Definition}
\newtheorem{example}[subsection]{Example}
\newtheorem{exercise}[subsection]{Exercise}
\newtheorem{situation}[subsection]{Situation}

\theoremstyle{remark}
\newtheorem{remark}[subsection]{Remark}
\newtheorem{remarks}[subsection]{Remarks}

\numberwithin{equation}{subsection}

% Macros
%
\def\lim{\mathop{\mathrm{lim}}\nolimits}
\def\colim{\mathop{\mathrm{colim}}\nolimits}
\def\Spec{\mathop{\mathrm{Spec}}}
\def\Hom{\mathop{\mathrm{Hom}}\nolimits}
\def\Ext{\mathop{\mathrm{Ext}}\nolimits}
\def\SheafHom{\mathop{\mathcal{H}\!\mathit{om}}\nolimits}
\def\SheafExt{\mathop{\mathcal{E}\!\mathit{xt}}\nolimits}
\def\Sch{\mathit{Sch}}
\def\Mor{\mathop{\mathrm{Mor}}\nolimits}
\def\Ob{\mathop{\mathrm{Ob}}\nolimits}
\def\Sh{\mathop{\mathit{Sh}}\nolimits}
\def\NL{\mathop{N\!L}\nolimits}
\def\CH{\mathop{\mathrm{CH}}\nolimits}
\def\proetale{{pro\text{-}\acute{e}tale}}
\def\etale{{\acute{e}tale}}
\def\QCoh{\mathit{QCoh}}
\def\Ker{\mathop{\mathrm{Ker}}}
\def\Im{\mathop{\mathrm{Im}}}
\def\Coker{\mathop{\mathrm{Coker}}}
\def\Coim{\mathop{\mathrm{Coim}}}

% Boxtimes
%
\DeclareMathSymbol{\boxtimes}{\mathbin}{AMSa}{"02}

%
% Macros for moduli stacks/spaces
%
\def\QCohstack{\mathcal{QC}\!\mathit{oh}}
\def\Cohstack{\mathcal{C}\!\mathit{oh}}
\def\Spacesstack{\mathcal{S}\!\mathit{paces}}
\def\Quotfunctor{\mathrm{Quot}}
\def\Hilbfunctor{\mathrm{Hilb}}
\def\Curvesstack{\mathcal{C}\!\mathit{urves}}
\def\Polarizedstack{\mathcal{P}\!\mathit{olarized}}
\def\Complexesstack{\mathcal{C}\!\mathit{omplexes}}
% \Pic is the operator that assigns to X its picard group, usage \Pic(X)
% \Picardstack_{X/B} denotes the Picard stack of X over B
% \Picardfunctor_{X/B} denotes the Picard functor of X over B
\def\Pic{\mathop{\mathrm{Pic}}\nolimits}
\def\Picardstack{\mathcal{P}\!\mathit{ic}}
\def\Picardfunctor{\mathrm{Pic}}
\def\Deformationcategory{\mathcal{D}\!\mathit{ef}}
\usepackage{cleveref}
\crefname{originallemma}{original lemma}{original lemmas}
\Crefname{originallemma}{Original lemma}{Original lemmas}
\makeatletter
\newlabel{lemma-NAK}{{00DV}{source}{}{sourcepointer-1}{}}
\newlabel{lemma-NAK@cref}{{[originallemma][0][]00DV}{[1][0][]source}{}{}{}}
\makeatother
\crefname{originallemma}{original lemma}{original lemmas}
\Crefname{originallemma}{Original lemma}{Original lemmas}
\makeatletter
\newlabel{lemma-ff}{{00HP}{source}{}{sourcepointer-2}{}}
\newlabel{lemma-ff@cref}{{[originallemma][0][]00HP}{[1][0][]source}{}{}{}}
\makeatother
\crefname{originallemma}{original lemma}{original lemmas}
\Crefname{originallemma}{Original lemma}{Original lemmas}
\makeatletter
\newlabel{lemma-local-criterion-flatness}{{00MK}{source}{}{sourcepointer-3}{}}
\newlabel{lemma-local-criterion-flatness@cref}{{[originallemma][0][]00MK}{[1][0][]source}{}{}{}}
\makeatother
\crefname{originallemma}{original lemma}{original lemmas}
\Crefname{originallemma}{Original lemma}{Original lemmas}
\makeatletter
\newlabel{lemma-variant-local-criterion-flatness}{{00ML}{source}{}{sourcepointer-4}{}}
\newlabel{lemma-variant-local-criterion-flatness@cref}{{[originallemma][0][]00ML}{[1][0][]source}{}{}{}}
\makeatother
\crefname{originallemma}{original lemma}{original lemmas}
\Crefname{originallemma}{Original lemma}{Original lemmas}
\makeatletter
\newlabel{remark-Tor-ring-mod-ideal}{{00M6}{source}{}{sourcepointer-5}{}}
\newlabel{remark-Tor-ring-mod-ideal@cref}{{[originallemma][0][]00M6}{[1][0][]source}{}{}{}}
\makeatother
\title{The fibre criterion for flatness in the Noetherian case}
\author{The Stacks Project Authors}
\date{stacks-project:algebra.tex:local-snapshot-20261012; GNU FDL 1.2 or later}
\begin{document}
\maketitle
\noindent\textit{Editorial note.} All delivered mathematical text is the author's own, copied from the pinned Stacks Project chapter file algebra.tex; no statement or proof step was rewritten, reordered or removed.\par\medskip
\noindent\textit{Editorial note.} Preamble provenance, as required for this source: the frozen snapshot stages the chapter's own preamble file beside it (preamble.tex, 217 lines, 8116 bytes, the same shared Stacks-wide preamble the published chow and etale-cohomology cards use); the first line of the pinned chapter file is \textbackslash{}input\{preamble\}, so no preamble text lives inside the chapter itself. The card takes that file's lines 8-18 (the verbatim package and the three \textbackslash{}comment-based environments reference/slogan/history at lines 9-12, section 13.2, together with the Xy-pic setup) and its lines 154 to the runtime end of the file (the native theorem declarations and mathematical macros). Not one character was changed, no package was added, no macro was defined, and the excerpt serves only as preamble, never as a span coordinate.\par\medskip
\noindent\textit{Editorial note.} The same-chapter pointers are located at their \textbackslash{}label lines inside the pinned chapter file; the cross-chapter pointers are located in the project's own tag list staged in the snapshot (sources/stacks-chow/tags.txt), at the verbatim lines carrying their canonical tags.\par\medskip
\noindent\textit{Editorial note.} The counted proof refers to same-chapter statements by label; they are recorded as navigation stubs pointed at their \textbackslash{}label lines, and the arguments consuming them stay inside the counted span.\par\medskip
\noindent\textit{Editorial note.} The counted spans contain no \textbackslash{}cite, so the card carries no bibliography.\par\medskip
\noindent\textit{Allowed standard prerequisites (editorial).}\begin{itemize}
\item The local criterion for flatness. Invoked at the end of the counted proof (line 24239). Converts the vanishing of a Tor group against the residue field into flatness.
\item Finitely generated modules over local rings and Nakayama. Used in the counted proof (lines 24214-24236). Propagates flatness from the fibre to the whole module.
\end{itemize}\medskip
\ifdefined\hypertarget\hypertarget{sourcepointer-1}{}\fi
\noindent\textit{Source reference.} Original lemma 00DV: sources/stacks-algebra/pinned-original-source.tex, source lines 3833--3833. Navigation pointer for \textbackslash{}ref\{lemma-NAK\} used inside the counted span (line 24224); the argument that consumes it is retained in the span.\par\medskip
\ifdefined\hypertarget\hypertarget{sourcepointer-2}{}\fi
\noindent\textit{Source reference.} Original lemma 00HP: sources/stacks-algebra/pinned-original-source.tex, source lines 9341--9341. Navigation pointer for \textbackslash{}ref\{lemma-ff\} used inside the counted span (line 24226); the argument that consumes it is retained in the span.\par\medskip
\ifdefined\hypertarget\hypertarget{sourcepointer-3}{}\fi
\noindent\textit{Source reference.} Original lemma 00MK: sources/stacks-algebra/pinned-original-source.tex, source lines 23800--23800. Navigation pointer for \textbackslash{}ref\{lemma-local-criterion-flatness\} used inside the counted span (line 24239); the argument that consumes it is retained in the span.\par\medskip
\ifdefined\hypertarget\hypertarget{sourcepointer-4}{}\fi
\noindent\textit{Source reference.} Original lemma 00ML: sources/stacks-algebra/pinned-original-source.tex, source lines 23968--23968. Navigation pointer for \textbackslash{}ref\{lemma-variant-local-criterion-flatness\} used inside the counted span (line 24222); the argument that consumes it is retained in the span.\par\medskip
\ifdefined\hypertarget\hypertarget{sourcepointer-5}{}\fi
\noindent\textit{Source reference.} Original lemma 00M6: sources/stacks-algebra/pinned-original-source.tex, source lines 18907--18907. Navigation pointer for \textbackslash{}ref\{remark-Tor-ring-mod-ideal\} used inside the counted span (line 24221); the argument that consumes it is retained in the span.\par\medskip
\begin{lemma}[Crit\`ere de platitude par fibres; Noetherian case]
\label{lemma-criterion-flatness-fibre-Noetherian}
Let $R$, $S$, $S'$ be Noetherian local rings and let $R \to S \to S'$
be local ring homomorphisms. Let $\mathfrak m \subset R$ be the
maximal ideal. Let $M$ be an $S'$-module. Assume
\begin{enumerate}
\item The module $M$ is finite over $S'$.
\item The module $M$ is not zero.
\item The module $M/\mathfrak m M$
is a flat $S/\mathfrak m S$-module.
\item The module $M$ is a flat $R$-module.
\end{enumerate}
Then $S$ is flat over $R$ and $M$ is a flat $S$-module.
\end{lemma}

\begin{proof}
Set $I = \mathfrak mS \subset S$. Then we see that $M/IM$ is a flat
$S/I$-module because of (3). Since
$\mathfrak m \otimes_R S' \to I \otimes_S S'$ is surjective we see
that also $\mathfrak m \otimes_R M \to I \otimes_S M$ is surjective.
Consider
$$
\mathfrak m \otimes_R M \to I \otimes_S M \to M.
$$
As $M$ is flat over $R$ the composition is injective
and so both arrows are injective.
In particular $\text{Tor}_1^S(S/I, M) = 0$ see
Remark \ref{remark-Tor-ring-mod-ideal}. By
Lemma \ref{lemma-variant-local-criterion-flatness} we conclude
that $M$ is flat over $S$. Note that since $M/\mathfrak m_{S'}M$
is not zero by Nakayama's Lemma \ref{lemma-NAK}
we see that actually $M$ is faithfully flat over $S$ by
Lemma \ref{lemma-ff} (since it forces $M/\mathfrak m_SM \not = 0$).

\medskip\noindent
Consider the exact sequence
$0 \to \mathfrak m \to R \to \kappa \to 0$.
This gives an exact sequence
$0 \to \text{Tor}_1^R(\kappa, S) \to \mathfrak m \otimes_R S \to I \to 0$.
Since $M$ is flat over $S$ this gives an exact sequence
$0 \to \text{Tor}_1^R(\kappa, S)\otimes_S M \to
\mathfrak m \otimes_R M \to I \otimes_S M \to 0$.
By the above this implies that $\text{Tor}_1^R(\kappa, S)\otimes_S M = 0$.
Since $M$ is faithfully flat over $S$ this implies that
$\text{Tor}_1^R(\kappa, S) = 0$ and we conclude that
$S$ is flat over $R$ by Lemma \ref{lemma-local-criterion-flatness}.
\end{proof}

\end{document}
